\chapter{Background and Literature Review}
\label{ch:background}

\section{Distributed Systems}
\label{sec:bg_distributed_systems}

\subsection{Definition and Characteristics}
\label{subsec:bg_ds_definition}

Distributed systems are collections of independent computers that appear to users as a
single coherent system \cite{tanenbaum2007distributed}. Their autonomous components
communicate via message passing and may span diverse hardware, virtual machines, and
geographic locations. The foundational characteristics of such architectures are:

\begin{itemize}
    \item \textbf{Scalability}: The capacity to dynamically adapt to an expanding
    workload by adding distributed nodes (horizontal scaling) or enhancing existing
    node resources (vertical scaling) \cite{gorton2022foundations, kleppmann2017designing}.

    \item \textbf{Fault Tolerance}: The ability to sustain operational correctness
    despite component, node, or network failures, realized through structural redundancy
    and state replication \cite{kleppmann2017designing, tanenbaum2007distributed}.

    \item \textbf{Transparency}: Abstraction of the distributed topology from
    applications and users, encompassing location transparency (hiding resource
    placement) and access transparency (uniform access to local and remote assets)
    \cite{tanenbaum2007distributed, coulouris2012distributed}.

    \item \textbf{Concurrency}: Coordinated simultaneous access to shared resources
    by multiple independent processes, enforced through locking mechanisms or
    distributed consensus protocols \cite{coulouris2012distributed, kleppmann2017designing}.

    \item \textbf{Heterogeneity}: Seamless integration across diverse hardware
    architectures, operating environments, and network infrastructures through
    standardized communication protocols \cite{coulouris2012distributed, tanenbaum2007distributed}.
\end{itemize}

\subsection{The CAP Theorem}
\label{subsec:bg_cap}

The CAP Theorem, conjectured by Brewer \cite{brewer2000towards} and formally proven by
Gilbert and Lynch \cite{gilbert2002brewer}, states that any distributed data store can
simultaneously guarantee at most two of the following three properties:

\begin{itemize}
    \item \textbf{Consistency (C)}: Every read reflects the most recent write, making
    the distributed topology appear as a single logical entity to the client.

    \item \textbf{Availability (A)}: Every non-failing node responds to requests in a
    reasonable timeframe, without guaranteeing that the response reflects the most
    recent mutation.

    \item \textbf{Partition Tolerance (P)}: The system sustains continuous operation
    despite the loss or arbitrary delay of messages between network nodes.
\end{itemize}

Because network partitions are unavoidable in practice, architects must choose between
CP systems (prioritizing consistency over availability) and AP systems (prioritizing
availability while accepting transient data divergence) \cite{gilbert2002brewer,
kleppmann2017designing}.

\subsection{Common Failure Modes}
\label{subsec:bg_failure_modes}

Distributed failures manifest across several fault domains \cite{coulouris2011distributed}:

\begin{itemize}
    \item \textbf{Network Failures}: Partitions, packet loss, latency spikes, and DNS
    resolution failures that disrupt inter-service communication.

    \item \textbf{Node Failures}: Hardware degradation, abrupt crash failures, and
    crash-recovery loops where a node continuously restarts, potentially losing
    volatile state.

    \item \textbf{Software and Service Failures}: Application bugs, deadlocks,
    misconfigurations, and resource exhaustion (CPU, memory, or thread pool starvation).

    \item \textbf{Data and Consistency Issues}: Stale reads due to replication lag,
    split-brain scenarios during partitions, and Byzantine failures where nodes exhibit
    arbitrary or malicious behavior \cite{kleppmann2017designing, raynal2018fault}.
\end{itemize}

\section{Microservices Architecture}
\label{sec:bg_microservices}

\subsection{Definition and Comparison with Monolithic Architecture}
\label{subsec:bg_ms_definition}

Microservices decompose applications into suites of small, autonomous services, each
encapsulating a singular business capability with its own codebase and persistence
layer \cite{newman2015building}. Services communicate exclusively via well-defined
network APIs, enabling independent deployment and scaling \cite{dragoni2017microservices}.
Compared to monolithic architectures---which co-locate all logic into a single
deployable artifact---microservices offer superior fault isolation, independent
horizontal scalability, and polyglot technology diversity, at the cost of significantly
increased operational complexity in monitoring, data consistency, and network overhead.

\begin{table}[H]
    \centering
    \begin{tabular}{|p{0.25\textwidth}|p{0.32\textwidth}|p{0.32\textwidth}|}
        \hline
        \textbf{Aspect} & \textbf{Microservices} & \textbf{Monolithic} \\
        \hline
        Structure & Small, independent services & Single, tightly integrated codebase \\
        \hline
        Deployment & Independent per service & Deployed as one unit \\
        \hline
        Scalability & Per-service & Entire application uniformly \\
        \hline
        Fault Tolerance & High (failure isolation) & Low (single point of failure) \\
        \hline
        Complexity & Operationally complex & Structurally simple \\
        \hline
    \end{tabular}
    \caption{Comparison of Microservices and Monolithic Architectures}
    \label{tab:ms_vs_monolith}
\end{table}

\subsection{Containerization and Orchestration}
\label{subsec:bg_ms_deployment}

Microservice deployment relies on containerization (primarily Docker
\cite{merkel2014docker}), which isolates applications and their dependencies into
lightweight, immutable environments. Orchestration platforms such as Kubernetes
\cite{burns2016borg} automate the lifecycle of containerized workloads, providing
automated rollouts, service discovery, internal load balancing, and self-healing
mechanisms.

\section{Observability in Distributed Systems}
\label{sec:bg_observability}

\subsection{The Three Pillars: Metrics, Logs, Traces}
\label{subsec:bg_pillars}

Observability is defined by a system's capacity to be understood through its external
outputs \cite{kalloniatis2020microservices_observability}. It is realized through three
interconnected telemetry pillars:

\begin{itemize}
    \item \textbf{Metrics}: Aggregated numerical health indicators (e.g., CPU
    utilization, latency percentiles), collected by tools such as Prometheus
    \cite{prometheus2016} and visualized in Grafana \cite{grafana2015}.

    \item \textbf{Logs}: Immutable, timestamped records of discrete computational
    events, aggregated via the ELK Stack (Elasticsearch, Logstash, Kibana)
    \cite{elk2013}.

    \item \textbf{Traces}: Reconstructions of the request lifecycle across service
    boundaries, enabling engineers to pinpoint bottlenecks \cite{sigelman2010dapper},
    implemented by tools such as Jaeger \cite{jaeger2017} and Zipkin \cite{zipkin2012}.
\end{itemize}

True observability links a high-level symptom (e.g., a metric anomaly indicating
increased latency) to its source pathway (e.g., a trace revealing a downstream
timeout), and finally to the granular cause (e.g., a log exposing a database
connection failure).

\subsection{Anomaly Detection Approaches}
\label{subsec:bg_anomaly_overview}

Automated anomaly detection spans a wide methodological spectrum:

\begin{enumerate}
    \item \textbf{Rule-Based Thresholds}: Static operational limits that are
    computationally trivial but prone to high false-positive rates under organic
    traffic spikes.

    \item \textbf{Statistical Methods}: Moving averages, rolling standard deviations,
    and z-scores that establish dynamic baselines, though they struggle with complex
    multivariate anomalies.

    \item \textbf{Predictive Analytics}: Leveraging historical seasonality to forecast
    expected metric behavior.

    \item \textbf{Machine Learning}: Unsupervised models (e.g., clustering) and deep
    learning (e.g., Autoencoders, LSTMs \cite{malhotra2015lstm}) capable of detecting
    minute, multidimensional deviations across vast data streams \cite{lokiny2026monitoring}.
\end{enumerate}

\section{Related Work}
\label{sec:bg_related_work}

\subsection{Anomaly Detection in Microservices}
\label{subsec:rw_anomaly}

Contemporary research emphasizes unsupervised and semi-supervised machine learning for
microservice environments, given the rarity of labeled fault data \cite{lin2018aiops}.
Statistical bounds have largely been superseded by sequential neural models,
particularly LSTM-Autoencoders, which exhibit superior competence in modeling the
temporal dependencies of service metrics.

\subsection{Federated Learning for System Monitoring}
\label{subsec:rw_federated}

Centralizing telemetry from geographically dispersed edge nodes incurs prohibitive
bandwidth costs and introduces data privacy concerns. Federated Learning (FL)
\cite{mcmahan2017communication} addresses this by enabling edge nodes to train models
locally and share only gradient updates with a central aggregator. Integrating
Differential Privacy \cite{dwork2006differential} into this pipeline \cite{abadi2016deep}
ensures that aggregated weights cannot be reverse-engineered to expose sensitive local
telemetry.

\subsection{Root Cause Analysis Systems}
\label{subsec:rw_rca}

Prominent RCA systems leverage topological service graphs derived from distributed
traces \cite{sigelman2010dapper, zhang2019rootcause}. Systems such as MicroRank
construct causal dependency graphs and employ random-walk algorithms akin to PageRank
\cite{page1999pagerank} to quantitatively score and rank probable fault origins amid
cascading failure alerts \cite{li2018microrank}.

\subsection{Adaptive Thresholding and Reinforcement Learning}
\label{subsec:rw_adaptive}

Static detection thresholds degrade rapidly in volatile cloud environments.
Reinforcement learning---specifically Q-learning \cite{watkins1992qlearning}---models
the thresholding problem as a Markov Decision Process, enabling agents to autonomously
adapt sensitivity profiles in real-time by learning from historical alert outcomes
\cite{sutton2018reinforcement, liu2020rlthreshold}.
